\section{结论}
\label{sec:conclusion}

我们引入二聚化 XXZ 链，用归一化反射标记 $\tilde{Z}_{\mathcal R}$ 画出
$L=8$ 有限尺寸相图，分辨平庸、拓扑与 AFM 三相。能隙标度确认有限尺寸效应
只动相边界。与需全层析的标记不同，我们的真机协议完全建立在可直接测量的
关联子上：保对称 orbit 拟设从匹配初态出发制备三相，优质比特筛选选出最优
真机子图，读出缓解展开原始计数。
该协议刻意模块化——更好的拟设、更丰富的筛选指标、关联读出模型都可各自升级
而不动比较框架——为在更大格点上追踪拓扑标记提供模块化起点。
本研究的局限有三：相划分全系于 $L=8$、尺寸外推只有能隙标度一路；
读出缓解假设各比特误差不关联；真机-模拟机趋势比较待 D08 图，
在此之前真机的比较一半是预留的协议而非结果。
